\documentclass[11pt,a4paper]{article}
\usepackage{fontspec}
\usepackage{kotex}
\usepackage{geometry}
\usepackage{graphicx}
\usepackage{xcolor}
\usepackage{tcolorbox}
\usepackage{booktabs}
\usepackage{array}
\usepackage{longtable}
\usepackage{colortbl}
\usepackage{fancyhdr}
\usepackage{titlesec}
\usepackage{hyperref}
\usepackage{emoji}
\usepackage{amsmath}
\usepackage{amssymb}
\usepackage{enumitem}

\geometry{left=25mm, right=25mm, top=30mm, bottom=30mm}

\setmainfont{Noto Sans CJK KR}
\setsansfont{Noto Sans CJK KR}
\setmonofont{Noto Sans CJK KR}
\setemojifont{Noto Color Emoji}

\definecolor{primary}{HTML}{1976D2}
\definecolor{darkbrown}{HTML}{3E2723}
\definecolor{mediumbrown}{HTML}{5D4037}
\definecolor{lightbrown}{HTML}{8D6E63}
\definecolor{cream}{HTML}{FFF3E0}
\definecolor{lightgray}{HTML}{F5F5F5}
\definecolor{tableheader}{HTML}{E8E8E8}

\titleformat{\section}
  {\Large\bfseries\color{primary}}
  {\thesection.}{0.5em}{}
\titleformat{\subsection}
  {\large\bfseries\color{darkbrown}}
  {\thesubsection}{0.5em}{}
\titleformat{\subsubsection}
  {\normalsize\bfseries\color{mediumbrown}}
  {\thesubsubsection}{0.5em}{}

\renewcommand{\contentsname}{목차}
\renewcommand{\figurename}{그림}
\renewcommand{\tablename}{표}

\tcbset{
  tipbox/.style={colback=cream,colframe=primary,boxrule=1.5pt,arc=3mm,
    left=5mm,right=5mm,top=3mm,bottom=3mm,fonttitle=\bfseries\color{primary}},
  formulabox/.style={colback=lightgray,colframe=lightgray,boxrule=0pt,arc=2mm,
    left=5mm,right=5mm,top=3mm,bottom=3mm},
  summarybox/.style={colback=white,colframe=primary,boxrule=2pt,arc=4mm,
    left=5mm,right=5mm,top=5mm,bottom=5mm}
}

\hypersetup{colorlinks=true,linkcolor=primary,urlcolor=primary,bookmarks=true,unicode=true}


\pagestyle{fancy}
\fancyhf{}
\fancyhead[R]{\small\color{lightbrown}마이크로비트 Grove I2C 센서 시뮬레이터 사용설명서 v1.0.0.0}
\fancyfoot[C]{\thepage}
\renewcommand{\headrulewidth}{0.4pt}

\begin{document}

\begin{titlepage}
\centering
\vspace*{3cm}
{\Huge\bfseries\color{primary}\emoji{electric-plug} 마이크로비트 Grove I2C 센서 시뮬레이터}\\[0.5cm]
{\Huge\bfseries\color{primary} 사용설명서}
\\[1.5cm]
{\Large 시뮬레이터 버전: v9.2.0.0}
\\[2cm]
\begin{tcolorbox}[summarybox]
\centering
{\large 이 사용설명서는 마이크로비트 Grove I2C 센서 시뮬레이터의 전체 기능을 상세히 설명합니다.}
\end{tcolorbox}
\vfill
{\large 사물인터넷과 빅데이터}\\[0.3cm]
{\large 청주교육대학교 컴퓨터교육과}
\vspace{1cm}
\end{titlepage}

\section*{1. 개요}

\subsection*{1.1 목적}

이 문서는 \textbf{마이크로비트 Grove I2C 센서 시뮬레이터}의 전체 기능을 상세히 설명하는 사용설명서입니다. 시뮬레이터를 처음 사용하기 전에 읽거나, 특정 기능을 자세히 알고 싶을 때 참조합니다.


\subsection*{1.2 관련 문서}

\begin{center}
\renewcommand{\arraystretch}{1.4}
\begin{tabular}{|p{5.5cm}|p{7.3cm}|p{2.3cm}|}
\hline
\rowcolor{tableheader}
\textbf{문서} & \textbf{용도} & \textbf{깊이} \\
\hline
퀵가이드 & 소프트웨어 사용 \textbf{중} 빠른 참조 & \emoji{star} \\
\hline
\textbf{사용설명서 (본 문서)} & \textbf{사용 전/후 전체 기능 이해} & \textbf{\emoji{star}\emoji{star}\emoji{star}} \\
\hline
이론해설서 & 배경 개념·원리 학습 & \emoji{star}\emoji{star}\emoji{star}\emoji{star} \\
\hline
활용가이드 & 수업 적용 방법 & \emoji{star}\emoji{star}\emoji{star} \\
\hline
개발참조문서 & 코드 구조·기술 사양 & \emoji{star}\emoji{star}\emoji{star}\emoji{star}\emoji{star} \\
\hline
\end{tabular}
\end{center}

\subsection*{1.3 시스템 요구사항}

\begin{center}
\renewcommand{\arraystretch}{1.4}
\begin{tabular}{|p{2.6cm}|p{12.4cm}|}
\hline
\rowcolor{tableheader}
\textbf{항목} & \textbf{요구사항} \\
\hline
브라우저 & Chrome, Edge, Firefox (최신 버전) \\
\hline
화면 & 1024px 이상 권장 (반응형 지원) \\
\hline
인터넷 & 불필요 (단일 HTML, 오프라인 사용 가능) \\
\hline
설치 & 불필요 (HTML 파일 더블클릭으로 실행) \\
\hline
\end{tabular}
\end{center}


\section*{2. 화면 구성}

\begin{center}
\includegraphics[width=0.95\textwidth]{images/S4_GroveSensor_screen.png}
\end{center}


\section*{3. 기능 상세}

\subsection*{3.1 센서 탐색}

10개 외장 센서 + 6개 내장 센서 카드 확인


\subsection*{3.2 드래그 앤 드롭}

센서 카드를 Grove Shield 포트(1\textasciitilde{}6)에 연결


\subsection*{3.3 키워드 추천}

13가지 카테고리에서 관련 센서 하이라이트


\subsection*{3.4 검색 필터}

텍스트 검색으로 센서 이름·설명 필터링


\subsection*{3.5 측정 데이터}

연결된 센서 조합의 데이터 항목 실시간 표시


\subsection*{3.6 데이터 조합}

최대 848가지 센서 조합 탐색 가능


\subsection*{3.7 포트 재배치}

연결된 센서를 다른 포트로 드래그 이동


\subsection*{3.8 초기화}

전체 포트 초기화 (모든 센서 제거)



\section*{4. 조작 방법}

\subsection*{4.1 기본 워크플로우}

1. \textbf{센서 탐색}: 우측 센서 카드 확인 또는 추천 키워드 활용

2. \textbf{센서 연결}: 카드를 드래그하여 중앙 포트에 드롭

3. \textbf{데이터 확인}: 하단 측정 데이터 패널에서 항목 확인

4. \textbf{조합 변경}: 포트의 센서 드래그로 이동 또는 제거



\section*{5. 자주 묻는 질문 (FAQ)}

\textbf{Q: 센서를 6개 넘게 연결할 수 있나요?}


A: Grove Shield의 I2C 포트가 6개이므로 최대 6개까지 연결 가능합니다.


\textbf{Q: 같은 센서를 2개 연결하면?}


A: 같은 I2C 주소를 가진 센서는 1개만 연결 가능합니다. 시뮬레이터에서 중복 연결을 방지합니다.


\textbf{Q: 내장 센서는 왜 드래그할 수 없나요?}


A: 내장 센서는 마이크로비트에 고정되어 있어 항상 사용 가능합니다. 별도 연결이 필요 없습니다.



\section*{6. 참고자료}

\begin{center}
\renewcommand{\arraystretch}{1.4}
\begin{tabular}{|p{1.5cm}|p{11.0cm}|p{2.5cm}|}
\hline
\rowcolor{tableheader}
\textbf{\#} & \textbf{자료명} & \textbf{비고} \\
\hline
1 & 마이크로비트Grove센서시뮬레이터\_퀵가이드 & 빠른 참조 \\
\hline
2 & 마이크로비트Grove센서시뮬레이터\_개발참조 & 기술 사양 \\
\hline
3 & 마이크로비트Grove센서시뮬레이터\_활용가이드 & 수업 적용 \\
\hline
4 & 마이크로비트Grove센서시뮬레이터\_이론해설서 & 배경 이론 \\
\hline
\end{tabular}
\end{center}


\vfill
\begin{center}
\small\color{lightbrown}
\textit{마이크로비트 Grove I2C 센서 시뮬레이터 사용설명서}
\end{center}
\end{document}